% Claim proof file for row claim_3_24 -- "SigmaSeparationEquivalences".
% Auto-generated stub by scaffold/solving_current_row.py at row start. A
% manager agent dedicated to writing the TeX proof fills in the body below.
% The proof must:
%   - Stay close to the lecture notes' proof if one exists (always search
%     the LN first for a `\begin{proof}` block immediately following the
%     claim; copy / lightly edit when present).
%   - Otherwise use the LN paradigm and cite prior claims by their `ref`.
%   - Be self-contained enough that a mathematician can follow it without
%     re-reading the LN.
%   - Have no `sorry`, `omitted`, or "obvious" shortcuts.
%
% Compiles standalone via the `subfiles` package.

\documentclass[main]{subfiles}

\begin{document}
% ref: claim_3_24 (proof, with statement restated for readability)
% title: SigmaSeparationEquivalences
%
% Cross-references: see claim_<ref>_statement_<title>.tex in this folder
% for the corresponding statement.

% --- statement (restated) ---------------------------------------------------
\begin{Rem}
    \begin{enumerate}
        \item By Proposition \ref{prp:sigma_opens} we have that $A \isPerp_G B \given C$ is equivalent to either of the following:
  \begin{enumerate}
      \item every walk from a node in $A$ to a node in $J \cup B$ is $C$-$\sigma$-blocked by $C$;
      \item every path from a node in $A$ to a node in $J \cup B$ is $C$-$\sigma$-blocked by $C$.
  \end{enumerate}
    \item Proposition \ref{prp:sigma_opens} also shows that if $A \nisPerp_G B \given C$ holds then:
  \begin{enumerate}
      \item there exists a (shortest) $C$-$\sigma$-open path from a node in $A$ to a node in $J \cup B$;
      \item there exists a (shortest) $C$-$\sigma$-open walk from a node in $A$ to a node in $J \cup B$ such that all its colliders lie in $C$.
  \end{enumerate}
  \end{enumerate}
  In practice we usually check if every path is $C$-$\sigma$-blocked or not. This is because there are, in contrast to walks, only a finite number of paths in a (finite) graph. In proofs, though, it often is easier to make use of walks, since these can be concatenated into walks (while one cannot in general concatenate two paths and again obtain a path).
\end{Rem}

% --- proof ------------------------------------------------------------------
\begin{proof}
% TODO: write the proof body.
\end{proof}

\end{document}
